\chapter{现代卷积神经网络}

LeNet 之后，更深的卷积网依靠更多数据与 GPU 才超过手工特征。
AlexNet 证明学习特征可行；VGG 用重复块，NiN 用 $1\times 1$ 卷积，
GoogLeNet 用并行 Inception，BatchNorm 稳定深层分布，ResNet 用 $\bm{x}+F(\bm{x})$，DenseNet 用通道连结。

\section{深度卷积神经网络（AlexNet）}

小数据上卷积早已有效，但在更大真实图像上，神经网络曾长期不敌支持向量机等。
传统视觉流水线用手工特征再接线性分类器；卷积则直接从像素学习 $f$。

\subsection{学习表征}

2012 年前主导的是 SIFT、SURF、HOG 与词袋等固定特征。
另一路线主张特征本身由数据学习，并把若干层叠成深度 $f_L\circ\cdots\circ f_1$。
每一层输出仍是特征图：每个空间--通道格点上的数是该层局部模式分数。

\subsubsection{缺少的成分：数据}

深度模型参数多，需要大量有标签样本才优于凸方法。
早期公开数据往往只有数百至数千张低分辨率图。
ImageNet 提供约 $10^{6}$ 张图、$10^{3}$ 类，使过参数化卷积网的经验风险可被数据约束。
$10^{6}$ 是样本个数，$10^{3}$ 是类别个数，都不是像素强度。

\subsubsection{缺少的成分：硬件}

每轮训练要多次穿过线性代数层。GPU 原本优化 $4\times 4$ 矩阵--向量积，与卷积中的局部乘加同类。
相对 CPU，GPU 以高吞吐并行乘加降低 $T_{\mathrm{device}}$（器件上的计算时间），使更深 $f$ 可训练。

\subsection{AlexNet}

AlexNet 为八层：五层卷积、两层全连接隐层、一层全连接输出，并在 ImageNet 上显著超过手工特征。
与 LeNet 同构但更深更宽，且用 $\operatorname{ReLU}$ 替代 sigmoid。

\subsubsection{模型设计}

输入按 $224\times 224$ 计，这两个整数是输入图像的空间高宽。
第一层核 $11\times 11$（匹配更大目标），随后 $5\times 5$ 与 $3\times 3$，这些是核的空间尺寸。
在第一、第二、第五卷积后接 $3\times 3$、步幅 $2$ 的最大汇聚：输出每个条目是窗口内最强响应，形状由填充--步幅公式给出空间高宽。
通道数约为 LeNet 的十倍。末两层全连接各 $4096$ 维，该整数是隐向量长度；每个分量是该隐单元的激活。
原实现曾把 $\bm{\theta}$ 拆到两块 GPU；显存充足时可放在单设备上。

\subsubsection{激活函数}

$\operatorname{ReLU}(u)=\max(0,u)$ 无需指数。
输入 $u$ 是仿射打分；输出在 $u>0$ 时等于该打分，否则为 $0$，表示“该单元未激活”。
sigmoid 在饱和区 $\sigma'(u)\approx 0$，该导数是损失经该点回传的局部斜率，接近 $0$ 表示参数几乎不被更新；
$\operatorname{ReLU}$ 在正半轴导数恒为 $1$，对初始化更不敏感。

\subsubsection{容量控制和预处理}

全连接层用暂退法约束复杂度；训练输入经翻转、裁切、变色等变换 $T$，
即在增广分布上最小化
\[
\mathbb{E}_{T}\bigl[\ell(f_{\bm{\theta}}(T(\bm{X})),y)\bigr].
\]
$\ell$ 是单样本损失（标量）；期望是对随机增广 $T$ 求平均。
该标量仍表示“增广后的预测与标签有多不符”，与未经增广时的损失同量纲。
用 $224\times 224$ 单通道张量穿过各层，可核验与架构表一致的 $\operatorname{shape}$：
表中的整数是批量、通道、高、宽这些轴长度，不是激活值。

\subsection{读取数据集}

原论文在 ImageNet 上训练。为缩短演示，把 Fashion-MNIST 的 $28\times 28$ 升到 $224\times 224$，
以便套用同一套核与步幅，尽管这并非该数据上的最优预处理。
这里改变的是空间高宽这两个整数，像素强度被重采样，但每个条目仍表示该格的亮度或颜色。

\subsection{训练AlexNet}

与 LeNet 相比，深度、宽度与分辨率都更大，故采用更小学习率 $\eta$，
否则 $\bm{\theta}-\eta\nabla L$ 容易跨过合适区域。
$\eta$ 是步长（每单位梯度把参数挪多远）；$\nabla L$ 的每个分量仍是该参数增大时损失大约增加多少。

\subsection{小结}

AlexNet 把 LeNet 式的 $\operatorname{Conv}\circ\operatorname{ReLU}\circ\operatorname{Pool}$ 加深加密，
并用暂退与增广控制 $\mathcal{F}$ 的有效容量。
卷积条目是局部模式分数，汇聚条目是最强响应，最后全连接给出类别打分。

\subsection{练习}

比较增加轮数后与 LeNet 的 $L$ 下降；核验主要显存占用是否在 $4096$ 维全连接的 $\bm{W}$，
主要乘加是否在前几层大分辨率卷积。
$L$ 是标量平均损失；$4096$ 是隐宽度（轴长度）。

\section{使用块的网络（VGG）}

AlexNet 证明深度有效，但未给出可复用的设计单元。
VGG 把“卷积--激活--汇聚”抽象成块，用重复块堆出不同深度。

\subsection{VGG块}

经典序列是：填充以保持分辨率的卷积、非线性、汇聚。
VGG 块由若干 $3\times 3$、填充 $1$ 的卷积（各接 $\operatorname{ReLU}$）再接 $2\times 2$、步幅 $2$ 的最大汇聚组成。
填充 $1$ 时
\[
n-3+2\cdot 1+1=n,
\]
左边按填充--步幅公式计算输出高度（或宽度），右边的 $n$ 表示该空间尺寸与输入相同，单位是格点数，不是响应强度。
故块内空间尺寸不变，块末汇聚使高宽减半：减半的是位置个数。
记卷积层数为 $n$、输出通道为 $c$，则一块是映射
\[
\mathrm{VGG\text{-}block}(\mathsf{X};n,c)
=
\operatorname{MaxPool}_{2,2}\bigl((\operatorname{ReLU}\circ\operatorname{Conv}_{3\times 3})^{n}(\mathsf{X})\bigr).
\]
输入 $\mathsf{X}$ 经 $n$ 次“局部模式打分再 ReLU”，再在每个 $2\times 2$ 窗口取最强响应。
输出每个条目仍是（更深的）局部模式摘要；形状中通道数为 $c$，高宽约为输入之半。

\subsection{VGG网络}

网络仍分卷积段与全连接段。超参数 $\mathrm{conv\_arch}$ 给出每块的 $(n,c)$。
$n$ 是块内卷积层数，$c$ 是该块输出通道数，二者都是计数。
原始 VGG-11 含五块：前两块各一层卷积，后三块各两层，通道 $64,128,256,512,512$，
故八层卷积加三层全连接。后续块将 $c$ 翻倍直至 $512$。

\subsection{训练模型}

完整 VGG-11 计算量大于 AlexNet。演示中减小各块通道数，仍在 Fashion-MNIST 上按同一损失训练，学习率可略高。
损失仍是分类交叉熵：模型输出经 softmax 后第 $k$ 个分量是属于类 $k$ 的概率。

\subsection{小结}

VGG-11 由可复用块 $\mathrm{VGG\text{-}block}(\,\cdot\,;n,c)$ 串联；更深更窄的 $3\times 3$ 栈优于更浅更宽的大核。
块输出的数是局部模式摘要，块参数 $(n,c)$ 是层数与通道个数。

\subsection{练习}

打印形状时卷积层计数为 $8$ 而非 $11$，核验其余三层是否为全连接；
并把输入从 $224$ 改为 $96$ 后代入步幅公式，看每块输出高宽如何减半。
代入后得到的整数是各块特征图的空间高宽。

\section{网络中的网络（NiN）}

LeNet、AlexNet、VGG 都是卷积提取空间特征再接全连接。
NiN 在每个空间位置的通道上再放多层感知机，即 $1\times 1$ 卷积，而不在早期展平。

\subsection{NiN块}

四维张量轴为样本、通道、高、宽。把每个像素当作样本、通道当作特征，则逐位置全连接即 $1\times 1$ 卷积。
NiN 块为
\[
\mathrm{NiN\text{-}block}
=
(\operatorname{ReLU}\circ\operatorname{Conv}_{1\times 1})^{2}
\circ
(\operatorname{ReLU}\circ\operatorname{Conv}_{k\times k}),
\]
先用 $k\times k$ 提取局部模式分数，再用两个 $1\times 1$ 在每个像素的通道维上做非线性混合。
每个输出条目是该位置、该通道上经过逐像素 MLP 后的特征响应；空间高宽由 $k\times k$ 卷积的填充与步幅决定。

\subsection{NiN模型}

卷积窗口取 $11\times 11$、$5\times 5$、$3\times 3$，通道设置平行于 AlexNet。
每块后接 $3\times 3$、步幅 $2$ 的最大汇聚，输出每个条目是窗口内最强响应，空间高宽按公式缩小。
取消末端全连接：最后一块输出通道数等于类别数 $q$，再全局平均汇聚
\[
o_c
=
\frac{1}{hw}\sum_{i=1}^{h}\sum_{j=1}^{w}[\mathsf{H}]_{c,ij},
\qquad
c=1,\ldots,q.
\]
$H_{c,ij}$ 是类别通道 $c$ 在位置 $(i,j)$ 的响应；$o_c$ 是该通道在整张图上的平均响应，作为类 $c$ 的 logit（未归一化打分）。
$h,w$ 是该特征图的空间高宽（格点数），$1/(hw)$ 把总和变成平均，不改变“这是类 $c$ 的整体分数”这一含义。
参数量显著小于大全连接，但有时增加达到相同训练误差的时间。

\subsection{训练模型}

在 Fashion-MNIST 上的目标与 AlexNet、VGG 相同，只是 $f_{\bm{\theta}}$ 换成 NiN 块加全局平均汇聚。
$\mathrm{softmax}(\bm{o})$ 的分量才是各类概率。

\subsection{小结}

NiN 用 $k\times k$ 后接 $1\times 1$ 的块提供逐像素 MLP，并用全局平均汇聚代替全连接输出。
$o_c$ 是类 $c$ 在空间上的平均响应，不是空间尺寸。

\subsection{练习}

去掉块中一个 $1\times 1$ 层，核验逐像素非线性深度如何改变；
并估计参数量主要来自各 $\operatorname{Conv}_{k}$ 的 $c_o c_i k^{2}$。
该乘积是单个核的参数个数（计数）。

\section{含并行连结的网络（GoogLeNet）}

不同任务上“最佳”核宽并不唯一。Inception 在同一层并行使用多种窗口，再沿通道连结。

\subsection{Inception块}

四条路径并行，输出在通道维连结：
\begin{align*}
\mathsf{P}_1
&=
\operatorname{Conv}_{1\times 1}(\mathsf{X}),\\
\mathsf{P}_2
&=
\operatorname{Conv}_{3\times 3}\bigl(\operatorname{Conv}_{1\times 1}(\mathsf{X})\bigr),\\
\mathsf{P}_3
&=
\operatorname{Conv}_{5\times 5}\bigl(\operatorname{Conv}_{1\times 1}(\mathsf{X})\bigr),\\
\mathsf{P}_4
&=
\operatorname{Conv}_{1\times 1}\bigl(\operatorname{MaxPool}_{3\times 3}(\mathsf{X})\bigr),\\
\mathsf{H}
&=
[\mathsf{P}_1,\mathsf{P}_2,\mathsf{P}_3,\mathsf{P}_4].
\end{align*}
各 $\mathsf{P}_r$ 的每个空间--通道条目是该路径检出的局部模式分数（或对汇聚摘要再混合后的分数）；
连结 $\mathsf{H}$ 不产生新数值，只把四路已有特征在通道维拼在一起。
中间路径的 $1\times 1$ 先把通道数降为 $c'$，$c'$ 是降维后的通道个数，再做 $3\times 3$ 或 $5\times 5$，以降低
乘加次数 $c_o c_i k^{2} n_h' n_w'$，其中 $n_h',n_w'$ 是输出特征图的空间高宽。
各路径输出高宽相同（靠填充），通道数相加：相加的是通道个数，不是把分数加在一起。

\subsection{GoogLeNet模型}

九个 Inception 块与其间的最大汇聚堆叠，末端用全局平均汇聚而非大全连接。
第一模块：$64$ 通道、$7\times 7$ 卷积，$64$ 是通道个数，$7\times 7$ 是核的空间尺寸。
第二模块：$1\times 1$ 扩到 $64$ 通道再 $3\times 3$ 把通道乘三，对应 Inception 第二条路径。
其后模块串联 Inception 块，块间汇聚降维：降的是空间格点数。通道分配比由实验选定。

\subsection{训练模型}

输入先变成 $96\times 96$，再在 Fashion-MNIST 上最小化同一分类损失。
$96$ 是空间边长（格点数）；损失仍是交叉熵。

\subsection{小结}

Inception 把多尺度 $\operatorname{Conv}_{1,3,5}$ 与汇聚并行后连结，并用 $1\times 1$ 降维；GoogLeNet 将其重复九次。
连结后每个条目仍是某一路径的局部模式分数。

\subsection{练习}

比较删去部分路径后的 Inception 与残差块的关系；并核验通道分配改变时 $\mathsf{H}$ 的通道维
$c_1+c_2+c_3+c_4$。
该和是四条路径输出通道个数之和，决定特征图有几条通道。

\section{批量规范化}

深层网难以在短时间内收敛。批量规范化（batch normalization）用小批量统计量校正中间表示，
并可与残差块一起支持很深的 $f$。

\subsection{训练深层网络}

对小批量 $\mathcal{B}$ 中的向量，
\begin{align*}
\hat{\bm{\mu}}_{\mathcal{B}}
&=
\frac{1}{|\mathcal{B}|}\sum_{\bm{x}\in\mathcal{B}}\bm{x},\\
\hat{\bm{\sigma}}_{\mathcal{B}}^{2}
&=
\frac{1}{|\mathcal{B}|}\sum_{\bm{x}\in\mathcal{B}}
(\bm{x}-\hat{\bm{\mu}}_{\mathcal{B}})^{2}
+\epsilon,\\
\mathrm{BN}(\bm{x})
&=
\bm{\gamma}\odot
\frac{\bm{x}-\hat{\bm{\mu}}_{\mathcal{B}}}{\hat{\bm{\sigma}}_{\mathcal{B}}}
+
\bm{\beta}.
\end{align*}
$\hat{\bm{\mu}}_{\mathcal{B}}$ 是该批量上每个通道（或特征）的样本均值，与激活同量纲，表示这一批里该通道的平均响应。
$\hat{\bm{\sigma}}_{\mathcal{B}}^{2}$ 是对应的样本方差（加 $\epsilon>0$ 避免除零），量纲是激活的平方，表示该通道在这一批里起伏有多大。
中间量 $(\bm{x}-\hat{\bm{\mu}}_{\mathcal{B}})/\hat{\bm{\sigma}}_{\mathcal{B}}$ 是按通道标准化后的激活：在这一批内该通道近似均值 $0$、方差 $1$，无量纲。
再经 $\bm{\gamma},\bm{\beta}$ 做可学习仿射：$\gamma_c$ 恢复该通道的尺度，$\beta_c$ 恢复偏移，使层不必被钉死在标准正态。
标准化针对的是激活值，不是空间高宽那些整数。

\subsection{批量规范化层}

BN 依赖整个小批量，不能在公式中把批量大小当作 $1$ 而忽略。
全连接与卷积上的归约轴不同。

\subsubsection{全连接层}

放在仿射与激活之间：
\[
\bm{h}
=
\phi\bigl(\mathrm{BN}(\bm{W}\bm{x}+\bm{b})\bigr).
\]
$\bm{W}\bm{x}+\bm{b}$ 的每个分量是线性打分；BN 先把它在批量维上按特征标准化成均值 $0$、方差 $1$，再仿射，最后 $\phi$ 给出非线性激活。
$\bm{\gamma},\bm{\beta}$ 与特征同维。

\subsubsection{卷积层}

同样放在卷积之后、$\phi$ 之前。每个输出通道一套标量 $\gamma_c,\beta_c$。
若批量大小为 $m$、空间为 $p\times q$，则第 $c$ 通道在 $m\cdot p\cdot q$ 个元素上共用同一
$\hat{\mu}_c,\hat{\sigma}_c$。
$m,p,q$ 是样本数与空间高宽（计数）；$\hat{\mu}_c$ 是这 $m\cdot p\cdot q$ 个激活的平均，$\hat{\sigma}_c$ 是其标准差。
标准化后该通道在整批、全空间上均值约 $0$、方差约 $1$；每个空间位置上的数仍是（已校正尺度的）局部模式响应。

\subsubsection{预测过程中的批量规范化}

预测时不再使用当前微批的噪声统计，且可能 $|\mathcal{B}|=1$。
用训练期移动平均 $\bm{\mu},\bm{\sigma}$ 代替 $\hat{\bm{\mu}}_{\mathcal{B}},\hat{\bm{\sigma}}_{\mathcal{B}}$，
使 $\mathrm{BN}$ 成为确定映射。
$\bm{\mu},\bm{\sigma}$ 是训练中见到的通道均值与标准差的平滑估计，含义与批量矩相同，只是不再随当前批抖动。
因此训练与预测的计算公式不同。

\subsection{从零实现}

实现对应上面的 $\mathrm{BN}_{\bm{\gamma},\bm{\beta}}$：训练时用批量矩并更新移动平均，
预测时用移动平均。$\bm{\gamma},\bm{\beta}$ 是参数；特征数在从零实现中需预先给出，该数是通道或特征的个数。

\subsection{使用批量规范化层的 LeNet}

在每个卷积或全连接之后、激活之前插入 BN。
训练目标不变，但中间表示被校正为（仿射前）按通道均值 $0$、方差 $1$，通常允许更大的 $\eta$。
第一层 BN 学到的 $\bm{\gamma},\bm{\beta}$ 可直接读出：$\gamma_c$ 是该通道被允许的尺度，$\beta_c$ 是偏移。

\subsection{简明实现}

框架中的 BN 与从零公式同一映射，只是核在编译后端上计算。
同一超参数下应得到同构的 $f_{\bm{\theta}}$。
读出的标准化激活同样是按通道均值 $0$、方差 $1$ 后再仿射的结果。

\subsection{争议}

BN 常使优化更平滑，并有正则化副作用。
将其解释为减少“内部协变量偏移”并不严谨：该说法与标准协变量偏移不是同一对象，
也未构成对泛化的完整证明。

\subsection{小结}

训练时
\[
\mathrm{BN}(\bm{x})
=
\bm{\gamma}\odot(\bm{x}-\hat{\bm{\mu}}_{\mathcal{B}})/\hat{\bm{\sigma}}_{\mathcal{B}}+\bm{\beta};
\]
括号内是按通道标准化后的激活（批量内均值 $0$、方差 $1$），再由 $\bm{\gamma},\bm{\beta}$ 恢复尺度与偏移。
预测改用全局移动矩。全连接与卷积的归约轴不同，但标准化的都是激活值，不是空间尺寸。

\subsection{练习}

核验仿射层在 BN 前是否仍需要偏置 $\bm{b}$（BN 已含 $\bm{\beta}$）；
并比较有无 BN 时可用的 $\eta$ 上界。
$\bm{\beta}$ 已经是标准化后的加性偏移，与 $\bm{b}$ 角色重叠。

\section{残差网络（ResNet）}

加深时必须保证新层不会把已有函数类变差。残差连接把恒等映射作为默认分量。

\subsection{函数类}

在函数类 $\mathcal{F}$ 中选
\[
f^{\ast}_{\mathcal{F}}
:=
\operatorname*{argmin}_{f\in\mathcal{F}}
L(\bm{X},\bm{y},f).
\]
$L$ 是标量训练损失；$f^{\ast}_{\mathcal{F}}$ 是该类里使损失最小的那个映射。
若 $\mathcal{F}_1\subset\mathcal{F}_2$，则
\[
\min_{f\in\mathcal{F}_2}L
\le
\min_{f\in\mathcal{F}_1}L.
\]
两边都是最小可达损失（非负标量）；嵌套类使更深模型至少能表示更浅模型。构造上需让恒等容易实现。

\subsection{残差块}

设理想映射为 $f(\bm{x})$。直接拟合 $f(\bm{x})$ 改为拟合残差 $F(\bm{x}):=f(\bm{x})-\bm{x}$，则
\[
f(\bm{x})=\bm{x}+F(\bm{x}).
\]
输入 $\bm{x}$ 是上一层传来的信号（激活或特征图）；$F(\bm{x})$ 是残差分支算出的修正量，与 $\bm{x}$ 同形状、同量纲。
相加后每个条目是原信号加上残差修正：若某处不需要改，最优是 $F\approx 0$，输出就等于原来的 $\bm{x}$。
该和是特征响应，不是空间尺寸相加。
若最优为恒等 $f(\bm{x})=\bm{x}$，只需 $F\approx\bm{0}$。
形状相同用恒等捷径；通道或步幅改变时用 $1\times 1$ 投影 $\bm{W}_s$：
\[
\bm{y}
=
\phi\bigl(\bm{x}+F(\bm{x})\bigr)
\quad\text{or}\quad
\phi\bigl(\bm{W}_s\bm{x}+F(\bm{x})\bigr).
\]
$\bm{W}_s\bm{x}$ 把捷径上的原信号改到与 $F(\bm{x})$ 相同的通道与空间尺寸，再逐元素相加；
$\phi$ 通常是 ReLU，输出每个条目仍是“校正后的信号再经非线性”。
块内 $F$ 通常为两层 $\operatorname{Conv}\circ\mathrm{BN}\circ\operatorname{ReLU}$。

\subsection{ResNet模型}

前两层同 GoogLeNet 风格：$64$ 通道、$7\times 7$ 卷积（步幅 $2$），再 $3\times 3$ 最大汇聚（步幅 $2$），
每卷积后接 BN。$64$ 是通道个数；$7,3$ 是核或窗口边长；步幅 $2$ 使空间高宽约减半。
其后四个模块由残差块组成：第一模块不降空间尺寸，之后每模块在首块把通道翻倍、高宽减半。
演示中每模块两个残差块。末端全局平均汇聚加全连接，得到类别打分。

\subsection{训练模型}

在 Fashion-MNIST 上最小化同一分类损失；前向为嵌套的 $\bm{x}\mapsto\bm{x}+F(\bm{x})$。
每一层输出的每个数都是原特征加上该层学到的残差修正。

\subsection{小结}

残差块实现 $f(\bm{x})=\bm{x}+F(\bm{x})$，输出中每个条目是原信号加残差修正，
使更深 $\mathcal{F}$ 容易包含恒等，从而可训练很深的网。

\subsection{练习}

比较 Inception 多路径连结与 $\bm{x}+F(\bm{x})$ 相加的差别；
并说明即使 $\mathcal{F}$ 嵌套，仍要限制 $F$ 的复杂度以免过拟合。
连结拼的是通道，相加要求形状相同并把对应位置的响应加在一起。

\section{稠密连接网络（DenseNet）}

DenseNet 把残差的加法换成通道维连结，使第 $\ell$ 层看见之前所有层的特征。

\subsection{从ResNet到DenseNet}

泰勒展开
\[
f(x)=f(0)+f'(0)x+\frac{f''(0)}{2!}x^{2}+\frac{f'''(0)}{3!}x^{3}+\cdots
\]
把 $f$ 拆成不同阶项。$f(0)$ 是常数项，$f'(0)x$ 是一阶项，其余为高阶；每一项与 $f(x)$ 同量纲。
ResNet 的一阶切分是
\[
f(\bm{x})=\bm{x}+g(\bm{x}),
\]
即原信号加一份修正。
DenseNet 把“加上 $g(\bm{x})$”换成与此前所有映射的连结
\[
\bm{x}
\;\longrightarrow\;
\bigl[
\bm{x},\;
f_1(\bm{x}),\;
f_2([\bm{x},f_1(\bm{x})]),\;
f_3([\bm{x},f_1(\bm{x}),f_2([\bm{x},f_1(\bm{x})])]),\;
\ldots
\bigr].
\]
连结不把响应相加，只增加通道：后面的层同时读到原始特征与各层新算出的特征，每个条目仍是各自的模式响应。
信息以拼接而非求和跨层传递。

\subsection{稠密块体}

块内使用 $\mathrm{BN}\circ\operatorname{ReLU}\circ\operatorname{Conv}$。
第 $i$ 个卷积块输出 $k$ 个通道（增长率 growth rate），并与输入在通道维连结。
$k$ 是每次新增加的通道个数。
若输入通道为 $c_0$、块内 $m$ 个卷积块，则输出通道
\[
c_0+m k.
\]
该和是通道个数（轴长度），例如 $c_0=3$、$m=2$、$k=10$ 时输出 $3+20=23$ 条通道，不是把 $23$ 加到像素上。

\subsection{过渡层}

稠密块使通道单调增加。过渡层用 $1\times 1$ 卷积把通道压到 $c'$，再用步幅 $2$ 的平均汇聚减半高宽：
\[
\mathsf{X}
\;\mapsto\;
\operatorname{AvgPool}_{2,2}\bigl(\operatorname{Conv}_{1\times 1}(\mathsf{X})\bigr).
\]
$1\times 1$ 卷积每个输出条目是逐像素通道混合后的分数；$c'$ 是压缩后的通道个数。
平均汇聚每个条目是 $2\times 2$ 窗口的平均响应，空间高宽减半指格点数减半。
从而控制模型复杂度。

\subsection{DenseNet模型}

茎部与 ResNet 相同：卷积加最大汇聚。随后四个稠密块，演示中每块 $4$ 层、增长率 $32$，故每块增加 $128$ 通道。
$4$、$32$、$128$ 都是计数（层数、每层新增通道、块内总新增通道）。
块间用过渡层减半空间并减半通道。末端全局汇聚加全连接，得到类别打分。

\subsection{训练模型}

网络较深时把输入高宽从 $224$ 降到 $96$ 以减小
$c\,h\,w$ 带来的显存，再在 Fashion-MNIST 上训练。
$c,h,w$ 分别是通道数与空间高宽，其乘积正比于特征张量的元素个数，不是激活的数值大小。

\subsection{小结}

ResNet 用 $\bm{x}+F(\bm{x})$，输出是原信号加残差修正；DenseNet 用通道连结
$[\bm{x},f_1(\bm{x}),f_2([\bm{x},f_1(\bm{x})]),\ldots]$，每个通道上的数仍是该层的特征响应，并由过渡层控制维数。

\subsection{练习}

解释参数量相对 ResNet 更小（核看到已拼接特征，无需再学一份副本）；
并核验输入改为 $224\times 224$ 时连结张量的显存如何随层数线性增长。
$224$ 是空间边长；显存随通道个数 $c_0+mk$ 线性增加。
